% Part: propositional-logic
% Chapter: syntax-and-semantics
% Section: semantic-notions

\documentclass[../../../include/open-logic-section]{subfiles}

\begin{document}

\olfileid{pl}{syn}{sem}

\olsection{అర్థపర భావనలు}

కింది అర్థపర భావనలను నిర్వచిస్తాం:

\begin{defn}
\begin{enumerate}
\item ఏదో~$\pAssign{v}$కు $\pSat{v}{!A}$ అయితే
  !!^a{formula}~$!A$ \emph{సంతృప్తిపరచదగినది}; ఏ
  $\pAssign{v}$కూ $\pSat{v}{!A}$ కాకపోతే అది
  \emph{సంతృప్తిపరచలేనిది};
\item అన్ని !!{valuation}s~$\pAssign{v}$కు $\pSat{v}{!A}$ అయితే
  !!^a{formula}~$!A$ ఒక \emph{సర్వసత్యం};
\item !!^a{formula}~$!A$ సంతృప్తిపరచదగినదై, సర్వసత్యం కాకపోతే అది
  \emph{పరిస్థిత్యాధీనమైనది};
\item $\Gamma$ ఒక !!{formula}s సమితి అయితే, $\pSat{v}{\Gamma}$ అయ్యే
  ప్రతి !!{valuation}~$\pAssign{v}$కు $\pSat{v}{!A}$ అయినప్పుడూ,
  అప్పుడు మాత్రమే $\Gamma \Entails !A$ (“$\Gamma$ నుంచి $!A$
  అర్థపరంగా అనుగమిస్తుంది”).
\item $\Gamma$ ఒక !!{formula}s సమితి అయితే, ఏదో
  !!a{valuation}~$\pAssign{v}$కు $\pSat{v}{\Gamma}$ అయినప్పుడు
  $\Gamma$ \emph{సంతృప్తిపరచదగినది}; లేకపోతే $\Gamma$
  \emph{సంతృప్తిపరచలేనిది}.
\end{enumerate}
\end{defn}

\begin{prob}
 కింది నాలుగు !!{formula}sలో ప్రతి ఒక్కటీ (a)~సంతృప్తిపరచదగినదా,
 (b)~సర్వసత్యమా, (c)~పరిస్థిత్యాధీనమైనదా అని నిర్ణయించండి.
\begin{enumerate}
  \item \( ( \Obj p_0 \lif ( \lnot \Obj p_1 \lif \lnot \Obj p_0 ) ) \).
  \item \( ( ( \Obj p_0 \land \lnot \Obj p_1 ) \lif ( \lnot \Obj p_0 \land \Obj p_2 )) \liff ( ( \Obj p_2 \lif \Obj p_0 ) \lif ( \Obj p_0 \lif \Obj p_1 )) \).
  \item \( ( \Obj p_0 \liff \Obj p_1 ) \lif ( \Obj p_2 \liff \lnot \Obj p_1 ) \).
  \item \( (( \Obj p_0 \liff ( \lnot \Obj p_1 \land \Obj p_2 )) \lor ( \Obj p_2 \lif ( \Obj p_0 \liff \Obj p_1 ))) \).
\end{enumerate}
\end{prob}

\begin{prop}
\ollabel{prop:semanticalfacts}
\begin{enumerate}
\item $!A$ ఒక సర్వసత్యం అయినప్పుడూ, అప్పుడు మాత్రమే
  $\emptyset \Entails !A$;
\item $\Gamma \Entails !A$, $\Gamma \Entails !A \lif !B$ అయితే
  $\Gamma \Entails !B$;
\item $\Gamma$ సంతృప్తిపరచదగినదైతే, $\Gamma$ యొక్క ప్రతి పరిమిత
  ఉపసమితి కూడా సంతృప్తిపరచదగినదే;
\item \ollabel{def:monotonicity} ఏకదిశత: $\Gamma \subseteq \Delta$,
  $\Gamma \Entails !A$ అయితే $\Delta \Entails !A$ కూడా;
\item \ollabel{def:Cut} సంక్రమణశీలత: $\Gamma \Entails !A$,
  $\Delta \cup \{ !A\} \Entails !B$ అయితే
  $\Gamma \cup \Delta \Entails !B$.
\end{enumerate}
\end{prop}

\begin{proof}
సాధన.
\end{proof}

\begin{prob}
\olref[pl][syn][sem]{prop:semanticalfacts}ను నిరూపించండి.
\end{prob}

\begin{prop}\ollabel{prop:entails-unsat}
  $\Gamma \Entails !A$ అయినప్పుడూ, అప్పుడు మాత్రమే
  $\Gamma \cup \{\lnot !A\}$ సంతృప్తిపరచలేనిది.
\end{prop}

\begin{proof}
సాధన.
\end{proof}

\begin{prob}
\olref[pl][syn][sem]{prop:entails-unsat}ను నిరూపించండి.
\end{prob}

\begin{thm}[అర్థపర నిగమన సిద్ధాంతం]
  \ollabel{thm:sem-deduction} $\Gamma \Entails !A \lif !B$
  అయినప్పుడూ, అప్పుడు మాత్రమే $\Gamma \cup \{!A\} \Entails !B$.
\end{thm}

\begin{proof}
సాధన.
\end{proof}

\begin{prob}
\olref[pl][syn][sem]{thm:sem-deduction}ను నిరూపించండి.
\end{prob}

\end{document}
